% Written from audit/F_E1_tame_construction.md and independently accepted by root.
\subsection{A repaired category using boundary-tame moduli paths}

The two changes in the next construction are explicit. In positive
characteristic we use paths tame at the boundary of the stable moduli
stack. In characteristic zero we restrict the covering arrows to those
whose Galois closures have order prime to the chosen residue
characteristic. Full characteristic-zero moduli markings remain allowed.
The obstructions E.1 and E.2 still apply to the original unrestricted
formulation.

\begin{claim}{E.10}{Proved}
Let $k$ be algebraically closed of characteristic $p>0$, and let
$\mathcal M=\mathcal M_{g,[n],k}$ be the stack of hyperbolic curves
with unordered punctures. Give its proper stable compactification
$\overline{\mathcal M}$ the divisorial logarithmic structure of its
boundary $D$. For $(m,p)=1$, put
\[
 \mathcal R_m=\sqrt[m]{(\overline{\mathcal M},D)},\qquad
 \mathcal P^{\mathrm{bt}}(\mathcal M)
   =\varprojlim_{(m,p)=1}\Pi_1^{\mathrm{et}}(\mathcal R_m),
\]
using geometric base points in $\mathcal M$. This is the full
boundary-tame fundamental groupoid. It is invariant under algebraically
closed field extension. For a smooth marked model over a strictly
henselian complete DVR of mixed characteristic $(0,p)$, it admits a
surjective specialization from the full geometric generic moduli
fundamental group, and the corresponding maps of path torsors are
surjective. Finally, the canonical map
\[
 \Aut_k(C)\longrightarrow
 \mathcal P^{\mathrm{bt}}(\mathcal M)(C,C)
\]
is injective for every hyperbolic curve $C/k$ of that type.
\end{claim}
\begin{proof}
We first justify the stack comparison being used. If $\mathcal X$ is
a proper finite-type Deligne--Mumford stack over a field or a complete
noetherian local ring, its diagonal is finite. Theorem~2.7 of
\citet{F_EHKV2001} gives a finite surjective representable map
$Y\to\mathcal X$ from a scheme. The scheme $Y$ and the next two
terms of its \v Cech nerve are proper schemes; the latter are schemes
because the diagonal is finite. Effective descent for finite
\emph{\'etale} covers along this map follows from SGA~1 IX,
Theorem~4.12, after an \emph{\'etale} atlas. It identifies a cover
of $\mathcal X$ with a finite \emph{\'etale} cover of $Y$, a
descent isomorphism on $Y\times_{\mathcal X}Y$, and its cocycle
on the triple product. Apply SGA~1 X, Corollary~1.8 to all three
proper schemes for algebraically closed field extension, or X,
Theorem~2.1 for passage from a complete local base to its closed
fibre. Full faithfulness carries the isomorphism and the cocycle
uniquely. Thus both finite-\emph{\'etale} comparison equivalences
hold for $\mathcal X$ \citep{E_SGA1}. This argument makes no
prime-to-$p$ assumption on the original stack inertia.

For completeness, the associated geometric specialization is
surjective when $\mathcal X$ is smooth proper over a complete DVR
$R$ with algebraically closed residue field. A connected finite
\emph{\'etale} special cover extends by the equivalence just
proved. The total cover is connected and regular. A connected
regular noetherian DM stack is irreducible, as may be checked on
\emph{\'etale} charts, so its generic fibre is connected.
Repeating after every finite extension of $R$ proves geometric
connectedness of that generic fibre. This is the Galois-category
criterion for surjectivity.

Knudsen's Theorem~2.7 supplies proper smoothness and relative normal
crossings for all stable pointed types; take the quotient by the
constant permutation group for unordered punctures
\citep[p.~179]{BC_Knudsen}. Locally, where the boundary parameters
are $x_1,\ldots,x_r$, the finite root stack is
\[
 \left[\Spec
 A[t_1,\ldots,t_r]/(t_1^m-x_1,\ldots,t_r^m-x_r)
       \big/\mu_m^r\right].
\]
The logarithmic construction glues these charts even when the
boundary has self-intersections or its local branches are permuted
\citep[Propositions~3.4 and 3.10]{F_TalpoVistoli2018}.
The stack is proper and smooth DM when $m$ is invertible. Its
finite representable \emph{\'etale} covers restrict fully faithfully
to the open moduli stack. Conversely, SGA~1 XIII, Proposition~5.2
extends a boundary-tame cover uniquely after extracting suitable
prime-to-$p$ roots of the boundary parameters. A common multiple
works on finitely many charts, and uniqueness of normalization
gives descent. This proves the groupoid description above. The
Kummer-flat theorem about infinite root stacks is not needed.

Apply the proper-stack comparisons to every $\mathcal R_m$.
They prove field-extension invariance. Over $R$, they give
surjections at each root level. The inverse-limit map is still
surjective: its compact image projects onto every level of a
directed inverse system and therefore is dense and closed.
Restriction of covers from a normal root stack to its dense open
also induces a surjection on fundamental groups. Combining these
maps proves the assertion for the full generic moduli group.
The assertion for path torsors follows by choosing a path and
using the group assertion. All ambiguities are the usual choices
of base points and paths.

It remains to prove the last assertion, which will exclude a
gerbe obstruction in E.11. Order the punctures; the ordering cover
extends as a finite \emph{\'etale} cover of the stable quotient
stack, even when $p\mid n!$. In genus zero an automorphism fixing
all these marks is identity. In positive genus add a full Jacobi
level at an odd prime $\ell\ne p$, $\ell\geq3$; in genus one
use an ordered mark as origin. These data kill automorphisms.
For genus at least two, use the faithful action on the Jacobian
in Deligne--Mumford Theorem~1.13 and Lemma~1.14. In genus one
use the faithful Tate-module action of origin-preserving elliptic
automorphisms. A finite-order matrix in
$1+\ell M_{2g}(\Z_\ell)$ is identity: the first nonzero
$\ell$-adic term of its powers, using $\ell$ odd, rules out both
prime-to-$\ell$ and $\ell$-power torsion. Thus a finite-order
automorphism fixing the full level acts trivially on the Tate
module and then on the curve
\citep[Theorem~1.13, Lemma~1.14, and \S5.14]{BC_DM}.

These level covers are boundary tame. On a regular chart of
the stable compactification, the universal curve is already
split semistable after strict localization. The prime-to-$p$
form of the monodromy computation in B.5 factors through tame
Kummer inertia, as in the proof of C.2. On homology modulo
$\ell$, a boundary generator acts by $1+N_e$, with
$N_e(x)=\langle x,\delta_e\rangle\delta_e$ and $N_e^2=0$.
Consequently $(1+N_e)^\ell=1$. At an intersection the
vanishing cycles are disjoint, so $N_eN_f=0$ and these
transvections generate an elementary abelian $\ell$-group.
Wild $p$-inertia is trivial. The level cover is therefore
tame at every boundary branch; the same argument applies
in the marked genus-one case. We use the prime-to-residue-
characteristic monodromy theorem underlying B.5 here, not a
claim about the full characteristic-$p$ geometric group
\citep[Theorem~4.8, Lemma~5.4]{BC_HMTopics}.
The action of $\Aut(C)$ on the fibre of the combined ordering
and level cover is faithful. Since this finite cover occurs in
the defining Galois category, its action detects every nonidentity
curve automorphism in $\mathcal P^{\mathrm{bt}}$.
\end{proof}

\begin{claim}{E.11}{Proved}
There is a category $\mathcal T_k^{\mathrm{bt}}$ with all
hyperbolic $k$-curves as objects, the groupoids
$\mathcal P^{\mathrm{bt}}$ as its marking groupoids, and all
actual finite \emph{\'etale} maps with prime-to-$p$ Galois
closure included faithfully. Every arrow factors as a marking
followed by such a cover. Covers transport along markings,
compatibly with composition. Algebraically closed field
extension induces an equivalence.

Write $c(X)=2g_X-2+n_X$. If $d=c(X)/c(Y)$ is not a positive
integer, then $\Hom(X,Y)$ is empty. Otherwise, for representatives
$u:V_u\to Y$ of the finitely many isomorphism classes of covers
of degree $d$, source type equal to that of $X$, and prime-to-$p$
Galois closure, there is a natural description
\[
 \Hom_{\mathcal T_k^{\mathrm{bt}}}(X,Y)
  =\coprod_u D_u\backslash
  \mathcal P^{\mathrm{bt}}(\mathcal M_{\mathrm{type}(X)})(X,V_u),
 \qquad D_u=\Aut_Y(V_u).
\]
These Hom spaces are compact Hausdorff, composition is continuous,
and $\Aut_{\mathcal T_k^{\mathrm{bt}}}(X)
=\pi_1^{\mathrm{bt}}(\mathcal M_{\mathrm{type}(X)},X)$.
\end{claim}
\begin{proof}
We give the path-lifting argument, including its inertia issue.
Let $u:V\to Y$ have Galois-closure group $G$ of order prime to
$p$. Use the stable Hurwitz stack with its Galois closure retained,
marking every puncture of $Y$ and every point above it. This is
the stabilized compactification of Bertin--Romagny, rather than
the generally nonnormal space of naive admissible covers.
Proposition~5.3, Definition~6.8, Theorem~6.9, Proposition~6.14,
Theorem~6.20 and Definitions~6.21--6.22 of \citet{E_BR2007}
give its properness, its source and target maps, and the
intermediate quotient constructions. All ramification points
are marked, so the nodal stabilizers are cyclic and do not
exchange branches. Their Theorem~5.5 gives the local
discriminant parameters
\[
                      \tau_i=t_i^{e_i},\qquad e_i\mid |G|,
\]
with the remaining deformation parameters unchanged.

Pull back a target root stack with $e_i\mid m$, and normalize.
The local branches are
$t_i=\zeta_i r_i^{m/e_i}$, where $\tau_i=r_i^m$.
Each branch has the same deformation ring as the root-stack
base. A relative automorphism fixing a branch fixes that full
deformation ring and extends uniquely over its universal
deformation. Its relative automorphism group is therefore the
deck group of the generic smooth cover. The extra boundary
automorphisms which permute the branches do not remain in the
relative stabilizer of a chosen normalization branch.
The identical computation for a non-Galois cover is
Bertin--Romagny Lemma~6.23 and Proposition~6.24. Consequently the
normalized morphism is, \emph{\'etale} locally on the target,
a finite disjoint union of finite \emph{\'etale} gerbes.
Equivalently it is such a gerbe over a representable finite
\emph{\'etale} cover of the target root stack. No
representability assertion is made for the Hurwitz discriminant
map itself; the nonrepresentability warning in Proposition~6.14
is essential here.

Because every $e_i$ is prime to $p$, the target roots so obtained
are cofinal among the prime-to-$p$ roots of the Hurwitz boundary.
The source map also acts on these tame groupoids: it extends
to the stable compactifications, and every source node parameter
pulls back to a monomial in Hurwitz node parameters times a unit.
The pullback of a prime-to-$p$ Kummer cover is again such a
cover. All inverse images of target marks have been retained;
this is what makes the intermediate quotient curves stable
and avoids the unmarked quotient problem noted in the source.

A finite \emph{\'etale} gerbe with fibre group $D$ has the
fundamental-group sequence
\[
       D\longrightarrow\pi_1(\mathcal H,u)
         \longrightarrow\pi_1(\mathcal Q,q)\longrightarrow1.
\]
The image of $D$ is normal and finite; this is Noohi
Proposition~10.4 and Corollary~10.5
\citep{F_Noohi2004}. The same assertion over a stack base is
proved by \emph{\'etale} descent: neutralize locally, check
it for $BD$, and observe that finite covers with trivial
relative-inertia action descend uniquely to $\mathcal Q$.
In general the left map need not be injective. Pass to the
cofinal root system first. The limit sequence remains exact:
surjectivity follows from the finite-intersection property
for compact lift sets; an element in the limit kernel comes
from a single element of the finite group $D$, again by the
finite-intersection property. The source map on this limit
sends $D$ to the canonical subgroup
$D\subset\Aut(V)\subset\pi_1^{\mathrm{bt}}(\mathcal M,V)$,
which is injective by E.10. Hence the left map is injective in
this application. This supplies the missing check that a
nontrivial gerbe transgression cannot kill deck transformations.

Thus a target tame path and a starting cover determine a unique
endpoint isomorphism class of covers and admit a lift in the
Hurwitz tame groupoid. Two lifts to the same endpoint differ by
a deck element. Their source paths therefore agree modulo the
action of that deck group.

Define an arrow by a pair $(u,\alpha)$, where $u:V\to Y$ is
one of the permitted covers and
$\alpha:X\rightsquigarrow V$ is a boundary-tame marking.
Identify $(u,\alpha)$ with $(u',h\alpha)$ when
$h:V\xrightarrow{\sim}V'$ is an isomorphism over $Y$.
For $(u,\alpha):X\to Y$ and $(v,\beta):Y\to Z$, write
$v:W\to Z$, transport $u$ along $\beta$ to
$u':V'\to W$, and let $\gamma:V\rightsquigarrow V'$ be its
source marking. Set
\[
                  (v,\beta)(u,\alpha)
                     =(vu',\gamma\alpha).
\]
The deck-kernel assertion proves independence of the lift and
of the chosen representatives. For associativity use the
Hurwitz stack of a chain of covers. A composite has Galois
closure contained in the corresponding wreath product, whose
order is still prime to $p$. Its common Galois closure and
the intermediate subgroups transport the chain simultaneously.
Their marked quotient models are stable by the result already
checked. Both bracketings give the same composite-cover map
and source path on that stack. Identities are $(\mathrm{id},
\mathrm{id})$.

Actual covers map to $(u,\mathrm{id})$. If two such maps became
equal, an isomorphism of their sources would have trivial
canonical marking. E.10 makes that isomorphism identity, so the
inclusion is faithful. Factorization and cover transport have
been built into the construction.

Tame Riemann--Hurwitz gives $c(V)=d\,c(Y)$. The maximal
prime-to-$p$ curve group is topologically finitely generated,
so it has only finitely many subgroups of a given index.
This proves the stated finite indexing of the Hom formula.
A path torsor is profinite and $D_u$ is finite; its quotient
is compact Hausdorff. On each component, endpoint transport
depends on a finite quotient, and after the endpoint is fixed
composition is induced by continuous maps of Hurwitz path
torsors and group multiplication. Composition is therefore
continuous. A degree-one cover is an isomorphism, giving the
asserted automorphism group.

Under $k\subset K$, E.3 identifies fixed-target covers, their
source types, and their deck groups. E.10 identifies the marking
torsors. The Hom formula proves full faithfulness. The moduli
stack of each hyperbolic type is geometrically connected, and
any $K$-curve has a type represented over $k$. A marking
connects it to such a base-changed curve. This proves essential
surjectivity, without asserting that the individual curve
descends as a scheme.
\end{proof}

\begin{claim}{E.12}{Proved}
Fix $p$, and let $\mathcal T^{\mathrm{cov}(p')}$ be the
characteristic-zero covering category with its full profinite
moduli markings, retaining precisely the covering arrows whose
Galois closures have order prime to $p$. After choosing good
reduction base models for the types and paths between type
base points, there is a specialization functor
\[
                 \mathcal T^{\mathrm{cov}(p')}
                     \longrightarrow \mathcal T_k^{\mathrm{bt}}.
\]
It preserves type and is surjective on each Hom space between
the chosen type base objects. On automorphism groups it is
\[
 \wh{\operatorname{Mod}}_{g,[n]}
       \twoheadrightarrow\pi_1^{\mathrm{bt}}(\mathcal M_{g,[n],k}).
\]
On Hom spaces it preserves the finite cover-component index
set and is the disjoint union of the surjections
\[
 D_u\backslash\operatorname{Path}_{\mathrm{et}}(X,V_u)
    \longrightarrow
 D_{u_s}\backslash\operatorname{Path}_{\mathrm{bt}}(X_s,V_{u_s}).
\]
The deck groups correspond under specialization. With compatible
path identifications and a group map $q:\Gamma\twoheadrightarrow
\Gamma_s$, the fibre above the coset of $q(\gamma)$ consists of
the left-$D_u$ classes in $D_u\ker(q)\gamma$.
\end{claim}
\begin{proof}
Take $R=W(k)$. For every type, smoothness of the pointed moduli
stack lifts a chosen special curve to a smooth marked model
over $R$; use a smooth chart through its point and completeness,
or formal existence for the proper curve. Fix one such model
per type. The full subcategories on their geometric generic
fibres and their special fibres are equivalent to the respective
full categories, because markings connect all objects of one
type. Characteristic-zero base-extension invariance identifies
the generic category with the stated one. One may use a common
algebraically closed overfield; no embedding of an arbitrarily
large field into $\C$ is required.

Let $u:V\to Y$ be a permitted cover of a chosen generic target.
By E.3 it corresponds to a special cover with the same degree,
source type and deck group. It has a smooth marked model over
the chosen target model after a finite extension of $R$:
lift the corresponding special cover, identify its geometric
generic fibre with $u$, and descend that identification to a
finite extension. Uniqueness of the marked stable model makes
this model and its deck action compatible with further
extensions.

Represent an arrow by $(u,\alpha)$, with $\alpha$ now a full
characteristic-zero moduli marking. Both endpoints of $\alpha$
have models after the preceding finite extension. E.10 specializes
the marking through the proper root stacks. Define
\[
                      (u,\alpha)\longmapsto
                         (u_s,\operatorname{sp}(\alpha)).
\]
The proper finite-\emph{\'etale} comparisons are compatible
with extension of the DVR. The separated stable moduli stack
gives the unique extensions of generic model isomorphisms.
These two facts give independence from extensions, representatives
and model identifications. The normalized root Hurwitz source
and target maps commute with specialization. Applying that
compatibility to the common-closure chain used in E.11 proves
compatibility with composition; hence this is a functor.

Degree and the marked source type are preserved by E.3. For
automorphisms, E.10 gives the surjection from the full generic
moduli group; the characteristic-zero comparison identifies it
with the full profinite mapping class group. More explicitly
the map first passes to the inverse limit of generic root
groups with indices prime to $p$, and then specializes. It
retains the target's finite interior $p$-monodromy; it is not
the maximal prime-to-$p$ quotient of the moduli group.

Finally apply the normal forms of E.11 and the fixed-target
cover equivalence E.3. Every cover component corresponds to
exactly one special component, with its same deck group.
Specialization of path torsors is surjective by E.10. It
therefore induces the displayed surjection on each component.
After identifying the torsors with their groups, two classes
have equal images exactly when their representatives differ
by a deck element and an element of $\ker(q)$, giving the
stated fibres. No injectivity of $q$ or faithfulness of this
specialization functor is asserted. E.2 proves that its
restriction on covering arrows cannot simply be removed.
\end{proof}
